\documentclass[11pt]{article}

\usepackage[utf8]{inputenc}
\usepackage[T1]{fontenc}
\usepackage{lmodern}
\usepackage{amsmath}
\usepackage{amssymb}
\usepackage{amsthm}
\usepackage[hidelinks]{hyperref}
\usepackage{doi}

% Map fi/ffi/fl ligature glyphs to Unicode so PDF full-text extraction (Zenodo indexing) is correct.
\input{glyphtounicode}
\pdfgentounicode=1

\newtheorem{theorem}{Theorem}
\newtheorem{proposition}{Proposition}
\newtheorem{lemma}{Lemma}
\newtheorem{corollary}{Corollary}
\theoremstyle{definition}
\newtheorem{definition}{Definition}
\theoremstyle{remark}
\newtheorem{remark}{Remark}

\newcommand{\CV}{\operatorname{CV}}
\newcommand{\rl}{r_{\ell}}
\newcommand{\Cgen}{\mathbb{C}^{3}_{\mathrm{gen}}}
\newcommand{\SPi}{S_{\Pi}}

\title{The Charged-Lepton Koide Relation as a Unit-Dispersion Constraint}

\author{J. Beau\\
Independent Researcher, France\\
\texttt{jerome.beau@cosmochrony.org}}

\date{2026}

\begin{document}

  \maketitle

  \begin{abstract}
    We record a scale-free reformulation of the charged-lepton Koide relation.
    Writing the square-root mass vector $\rl = (\sqrt{m_e}, \sqrt{m_\mu}, \sqrt{m_\tau})$ with mean $\bar r$ and
    standard deviation $\sigma_r$, the Koide quotient is exactly
    $Q_\ell = (1 + \CV(\rl)^2)/3$ with $\CV(\rl) = \sigma_r/\bar r$, so that
    $Q_\ell = 2/3 \iff \CV(\rl) = 1 \iff \sigma_r = \bar r$: the square-root masses have unit relative dispersion.
    This isolates the intrinsic shape content of Koide as a single dimensionless number and separates it cleanly
    from the mass hierarchy, which is fixed independently by the placement of one generation near a node of the
    square-root profile.
    We record two structural consequences for the Cosmochrony fermionic sub-programme.
    First, the specific cascade ladder $(1, \tfrac12 + u, \tfrac12 - u)$ that carries the projected generation
    levels has a bounded dispersion, $\CV \le 1/\sqrt2$ and $Q \le 1/2$, so it cannot by itself realise the Koide
    value.
    Second, $\CV(\rl) = 1$ is the coefficient of variation of the maximum-entropy positive variable at fixed mean;
    whether the projection entropy $\SPi$ induces such a variational principle on the square-root mass observable
    is stated as a sharp, falsifiable open test.
    This note does not derive the masses: it fixes a target invariant for the fermionic-mass sector.
  \end{abstract}

  \noindent\textbf{Keywords:}
  Koide relation; charged-lepton masses; coefficient of variation; unit dispersion; square-root mass vector;
  projective cascade; maximum entropy; projection entropy; fermionic mass target.

  \newpage

  \tableofcontents

  \newpage

  \section{Position and purpose}
    \label{sec:position}

    The charged-lepton mass sector of the Cosmochrony fermionic sub-programme is not derived: the projected Yukawa
    line fixes a generation-level structure, but the mass values and the mixing remain downstream open
    data~\cite{Beau2026pyl, Beau2026pyo}.
    In that setting it is useful to state the phenomenological target in the sharpest possible form, so that a
    later mechanism can be tested against a single scale-free invariant rather than against the raw mass ratios.
    The Koide relation is the natural candidate: it is dimensionless, three-generational, and formulated in the
    square-root masses.
    This note recasts it as a unit coefficient-of-variation condition, separates the resulting dispersion
    statement from the hierarchy, records the obstruction it poses to the projected cascade ladder, and isolates a
    maximum-entropy reading as the open structural test.
    The purpose is reconnaissance and bookkeeping, not derivation: the result is a target invariant, not a mass
    formula.

  \section{Koide as a unit coefficient of variation}
    \label{sec:cv}

    Let $\rl = (\sqrt{m_e}, \sqrt{m_\mu}, \sqrt{m_\tau})$, and write $\mathbf 1 = (1,1,1)$.
    The Koide quotient is
    \begin{equation}
      \label{eq:koide}
      Q_\ell
      = \frac{m_e + m_\mu + m_\tau}{(\sqrt{m_e} + \sqrt{m_\mu} + \sqrt{m_\tau})^2}
      = \frac{\lVert \rl \rVert^2}{\langle \rl, \mathbf 1 \rangle^2}.
    \end{equation}

    \begin{proposition}[Koide as unit dispersion]
      \label{prop:cv}
      Let $\bar r = \tfrac13 \sum_g r_g$ and $\sigma_r^2 = \tfrac13 \sum_g (r_g - \bar r)^2$ be the mean and
      population variance of the three components of $\rl$, and $\CV(\rl) = \sigma_r/\bar r$.
      Then
      \begin{equation}
        \label{eq:cv-identity}
        Q_\ell = \frac{1 + \CV(\rl)^2}{3},
      \end{equation}
      and consequently
      \begin{equation}
        \label{eq:cv-unit}
        Q_\ell = \frac23 \quad\Longleftrightarrow\quad \CV(\rl) = 1 \quad\Longleftrightarrow\quad \sigma_r = \bar r.
      \end{equation}
    \end{proposition}

    \begin{proof}
      With three components, $\lVert \rl \rVert^2 = \sum_g r_g^2 = 3(\bar r^2 + \sigma_r^2)$ and
      $\langle \rl, \mathbf 1 \rangle = \sum_g r_g = 3 \bar r$, so $\langle \rl, \mathbf 1 \rangle^2 = 9 \bar r^2$.
      Substituting into~\eqref{eq:koide} gives
      $Q_\ell = 3(\bar r^2 + \sigma_r^2)/(9 \bar r^2) = (1 + \sigma_r^2/\bar r^2)/3 = (1 + \CV(\rl)^2)/3$, which
      is~\eqref{eq:cv-identity}.
      Setting this equal to $2/3$ gives $\CV(\rl)^2 = 1$, hence~\eqref{eq:cv-unit}.
    \end{proof}

    Numerically, with the pole-mass square roots $\rl = (0.715, 10.28, 42.15)$ (in $\mathrm{MeV}^{1/2}$), one finds
    $\bar r = 17.716$, $\sigma_r = 17.715$, hence $\CV(\rl) = 0.999991$ and $Q_\ell = 0.666661$.
    Propagating the mass uncertainties (dominated by $m_\tau = 1776.86 \pm 0.12\,\mathrm{MeV}$) gives
    $\CV(\rl) = 0.99999 \pm 0.00001$, so $\CV = 1$ (equivalently $Q_\ell = 2/3$) is consistent with the data at
    $0.9\sigma$: the charged-lepton square-root mass vector has unit relative dispersion at the $10^{-5}$ level,
    limited by $m_\tau$ --- consistent with exact, not established as exact.
    Equation~\eqref{eq:cv-unit} is the sharpest scale-free statement of the relation: rather than ``why $3477$'' or
    ``why $45^\circ$'', the content is that the standard deviation of the square-root masses equals their mean.
    The value $Q_\ell$ is monotone in $\CV$ across the admissible interval $\CV \in [0, \sqrt2]$: $\CV = 0$ is the
    degenerate limit ($Q = 1/3$), $\CV = \sqrt2$ is the one-mass-dominant limit ($Q = 1$), and $\CV = 1$ is the
    Koide midpoint ($Q = 2/3$).

    \begin{remark}[Robustness and scope]
      \label{rem:robustness}
      Three qualifications fix the status of the equality.
      (i) \emph{Precision.} $\CV = 1$ holds at the $10^{-5}$ level and is consistent with exact at $0.9\sigma$, the
      precision limited by $m_\tau$; the tau mass making $\CV = 1$ exact is $1776.97\,\mathrm{MeV}$, within the
      measured $1776.86 \pm 0.12$.
      It is a robust near-equality consistent with exact, not a theorem that it is exact.
      (ii) \emph{Renormalisation scale.} $\CV = 1$ is an infrared (pole-mass) statement, not an RG invariant: the
      running masses at $M_Z$ give $\CV \simeq 1.002$.
      It is, however, invariant under any \emph{flavour-universal} rescaling of $\sqrt m$ (both $\sigma_r$ and
      $\bar r$ scale together), so the leading flavour-blind QED running leaves it exactly fixed and the small
      drift to $M_Z$ is a purely flavour-dependent (threshold, scheme) effect.
      Here ``scale-free'' means dimensionless, not RG-invariant.
      (iii) \emph{Sector specificity.} $\CV = 1$ is specific to the charged leptons; the quark sectors sit away
      from it (indicatively $\CV \simeq 1.09$ down-type, $\CV \simeq 1.24$ up-type, both scheme- and
      scale-dependent), and the neutrino $\CV$ is undefined without the absolute mass scale.
      A mechanism forcing $\CV = 1$ must therefore act at the infrared scale and be specific to the charged-lepton
      sector.
    \end{remark}

  \section{Separation between dispersion and hierarchy}
    \label{sec:separation}

    Proposition~\ref{prop:cv} fixes only the relative dispersion of the square-root masses.
    It does not determine which generation is light.
    In the cyclic parametrisation
    \begin{equation}
      \label{eq:cyclic}
      r_g = \bar r \left( 1 + \sqrt2 \, \cos\!\left( \delta + \frac{2\pi g}{3} \right) \right),
      \qquad g = 0, 1, 2,
    \end{equation}
    the condition $Q_\ell = 2/3$ fixes the amplitude to $\sqrt2$ for every phase $\delta$ (the three equally spaced
    cosines have vanishing sum and mean square $1/2$, so $\sigma_r/\bar r = 1$ identically), while $\delta$ places
    the three generations relative to the nodes of the square-root profile.
    Because the amplitude exceeds unity, $1 + \sqrt2 \cos\phi$ has real zeros at $\cos\phi = -1/\sqrt2$, i.e.\ at
    $\phi = 135^\circ$ and $\phi = 225^\circ$.
    The large hierarchy arises when one generation lies close to such a node.
    For the observed charged-lepton pole masses, the electron phase lies about $2.27^\circ$ from the $135^\circ$
    node, so $\sqrt{m_e}$ is nearly annulled by interference.
    Equivalently, the small angular offset is
    \begin{equation}
      \label{eq:eps-e}
      \varepsilon_e \simeq 3.96 \times 10^{-2}\ \mathrm{rad}.
    \end{equation}
    Near the node the square-root profile is linear in the angular offset, so
    \begin{equation}
      \label{eq:nodal}
      \sqrt{m_e} \simeq \bar r\,\varepsilon_e, \qquad m_e \simeq \bar r^2\,\varepsilon_e^2 .
    \end{equation}
    The hierarchy is therefore a nodal suppression: a small phase displacement is converted into a quadratic
    suppression of the lightest mass.
    This statement does not identify $\varepsilon_e$ with a known coupling constant; it records only the robust
    near-node mechanism.
    The additive-to-multiplicative bridge is logarithmic: $\ln m_g = 2 \ln\lvert 1 + \sqrt2 \cos\phi_g \rvert$ is
    smooth in the phase except for a logarithmic divergence at the nodes, so
    $\ln(m_\tau/m_e) \simeq -2 \ln(\text{node distance}) + O(1)$.
    The same fact explains the unequal steps: the $\mu/e$ step is steep because $e$ is near a node, while the
    $\tau/\mu$ step is mild because $\mu$ and $\tau$ are both off-node.

    \begin{remark}[Two independent targets]
      \label{rem:targets}
      The charged-lepton problem separates cleanly.
      The Koide relation is the dispersion target $\CV(\rl) = 1$; the concrete hierarchy is the phase statement
      that $\delta$ places one generation near a node.
      The first is scale-free and blind to which generation is light; the second is the residual order that
      selects the electron.
      A candidate mechanism should be tested against $\CV(\rl) = 1$ first, and against the node placement second.
      The two targets are not of equal difficulty.
      The dispersion target $\CV = 1$ is a single scale-free number with a plausible entropic reading
      (Section~\ref{sec:maxent}); the node target is a fine placement of one phase against a zero of the
      square-root profile, and near a node the mass ratio is exponentially sensitive to that placement, so
      ``why $\delta$ sets the electron at the node'' is very likely the harder problem and is where the genuine
      residue of the charged-lepton hierarchy lives.
      This note fixes the easier, scale-free target and defers the hard one; it does not claim the hard one is
      close.
    \end{remark}

    \paragraph{Physical picture: one three-phase wave, symmetric versus broken.}
      Equation~\eqref{eq:cyclic} invites a wave reading: the three generations are three equally spaced phases,
      $2\pi/3$ apart, sampled from one profile $1 + \sqrt2 \cos\phi$.
      In this language the two targets acquire a physical meaning.
      The dispersion $\CV = 1$ is the \emph{modulation depth} of the wave (amplitude $\sqrt2$), a single universal
      number blind to which generation is light.
      The hierarchy is \emph{interference}: the electron is small not because a factor crushes it, but because its
      phase falls near a dark fringe (a node), where the square-root amplitude nearly vanishes.
      The same three-phase skeleton then admits distinct projective regimes:
      \begin{center}
        \begin{tabular}{lll}
          \hline
          wave regime & projection & phenomenon \\
          \hline
          exactly equilateral phases & indiscernible classes & confined colour \\
          slightly imbalanced phases & discernible classes & massive generations \\
          one phase near a node & amplitude nearly annulled & electron smallness \\
          \hline
        \end{tabular}
      \end{center}
      In the exactly symmetric regime the three phases form a perfect $120^\circ$ triad with vanishing sum
      ($1 + \omega + \omega^2 = 0$): the three classes are indiscernible and none can be isolated~---~the structure
      of an unbroken colour triplet, whose non-observability as three separate charges is the confinement already
      carried by the non-injective colour projection of ENI~\cite{Beau2026l}.
      In the broken regime the phases are displaced off the perfect triad, the three amplitudes separate, and one
      reads three distinct masses.
      Confinement and the mass hierarchy would then be two faces of one three-phase wave: symmetric $\Rightarrow$
      colour (confined, never three values), broken $\Rightarrow$ generations (observed, three masses).
      This picture is heuristic.
      It explains the \emph{form} of the problem~---~why three, why the vanishing sum, why the electron is small
      (interference)~---~but not the \emph{values}: it fixes neither the modulation depth $\sqrt2$ nor the node
      phase $\delta$.
      Two cautions keep it honest.
      First, identifying the colour triad and the generation triad as one and the same wave is a non-standard
      step: in the Standard Model colour (a gauge triplet within one generation) and the generation index are
      independent.
      Second, the hard content stays exactly where Remark~\ref{rem:targets} places it, the fine node placement;
      and in the quantum-mechanical analogy the wave reading invokes, what is still missing is the operator or
      interaction that would \emph{force} the continuous phase profile to resolve into exactly three observed
      values with unit dispersion.
      In this picture the two open targets read as a \emph{form lock} and a \emph{placement lock}.
      The form lock is the modulation depth $A = \sqrt2$, and it is \emph{not} a condition separate from the Koide
      angle.
      Writing $\theta$ for the angle between $\rl$ and the democratic axis $(1,1,1)$, the three cosines
      $120^\circ$ apart have $\langle \cos^2 \rangle = 1/2$, so the oscillating part has effective size $A/\sqrt2$
      and
      \begin{equation}
        \label{eq:tan-theta}
        \begin{gathered}
          \tan\theta = \frac{\lVert \text{oscillating} \rVert}{\lVert \text{democratic} \rVert}
          = \frac{A}{\sqrt2} = \CV, \\[2pt]
          A = \sqrt2 \iff \CV = 1 \iff \tan\theta = 1 \iff \theta = 45^\circ .
        \end{gathered}
      \end{equation}
      The amplitude $\sqrt2$, the Koide angle $45^\circ$, and $\CV = 1$ are thus three writings of one content:
      the $\sqrt2$ is the natural three-phase factor (from $\langle \cos^2 \rangle = 1/2$), and the sole real
      condition is that the \emph{differentiating} (oscillating, generation-splitting) component of $\rl$ has
      exactly the same norm as its \emph{common} (democratic) component~---~a singlet--doublet equipartition.
      The form lock is therefore one question, not two: why does the projection place the
      generation-differentiation component exactly at parity with the democratic component?
      This is plausibly a global principle~---~entropy, saturation, equipartition, or a projective normalisation
      constraint.
      The phase $\delta$ (why the electron sits at the node) is the placement lock, more singular and fragile,
      since near a zero a small angular change produces a large mass ratio.
      The proposed bridge is stated cautiously, and this is its canonical form: \emph{not} that colour and
      generation are identical indices, but that they may be two projective regimes of the same three-phase
      carrier~---~exact equilateral phase gives confined colour, while broken phase gives observable generations
      and the mass hierarchy.

  \section{The cascade-ladder obstruction}
    \label{sec:ladder}

    The projected generation levels of the sub-programme are the eigenvalues of the squared projective residue on
    the gauge-singlet triplet $\Cgen$, in the normal form
    $(1, \tfrac12 + u, \tfrac12 - u)$ fixed by the Schur reduction with the even sector closed by Born--Infeld
    parity~\cite{Beau2026prs, Beau2026a34, Beau2026pyl}.
    This ladder has a bounded dispersion, independently of the split $u$.

    \begin{proposition}[Cascade-ladder bound]
      \label{prop:ladder}
      For the pinned ladder $(1, \tfrac12 + u, \tfrac12 - u)$ with $u \in [0, \tfrac12)$, read as either masses or
      square-root masses, the coefficient of variation satisfies $\CV \le 1/\sqrt2$, hence $Q \le 1/2$, with
      equality only at the singular endpoint $(1,1,0)$.
    \end{proposition}

    \begin{proof}
      The top eigenvalue is pinned at $1$ and the lower pair is centred at $1/2$.
      Within the admissible interval $[0,1]$ the dispersion is maximised by pushing the lower pair to the endpoint,
      i.e.\ $u \to 1/2$, giving $(1,1,0)$; there $\bar r = 2/3$, $\sigma_r^2 = 2/9$, so $\CV = 1/\sqrt2$ and, by
      Proposition~\ref{prop:cv}, $Q = 1/2$.
      Any larger value would require an eigenvalue above $1$ or a displacement of the lower-pair centre away from
      $1/2$, neither of which is available in the even cascade sector $\operatorname{diag}(1, \tfrac12, \tfrac12)$.
    \end{proof}

    The bound is a property of this specific pinned ladder, not of $\mathbb Z_2$-symmetric triples in general, and
    --- crucially --- it constrains the STATIC LEVEL structure, not the physical masses.
    The generation levels are the hermitian square of the projected Yukawa operator,
    $H_\Pi = \lambda_Y^2 \operatorname{diag}(1, \tfrac12 + u, \tfrac12 - u)$; their square roots have
    $\CV(\rl) \in [0.09, 0.37]$ at the dictionary value $u = 1/10$ (levels $1 : 3/5 : 2/5$), with $\lambda_Y$
    cancelling in the coefficient of variation.
    But the PHYSICAL charged-lepton mass carries, beyond these levels, the depth amplification of the projective
    cascade, $m_g = \lambda_Y\, w_g\, S_g\, A(n_g)$ with $A(n_g) = \exp(\beta^\ast n_g)$~\cite{Beau2026pyo}; its
    coefficient of variation is therefore set by the depth-vector $\{n_g\}$ and is NOT capped by the level ladder.
    Two consequences follow.
    First, the ladder bound $\CV \le 1/\sqrt2$ is a statement about the levels; presenting it as the obstruction to
    $\CV = 1$ would conflate the level structure with the mass.
    Second, within the corpus's present (depth) realisation of the cascade the physical $\CV$ is controlled by the
    same amplification depths that carry the hierarchy, so the forced $O(1)$ depths give neither $\CV = 1$ nor the
    hierarchy --- but this is a limitation of the depth realisation, not a foundational obstruction.
    Unit dispersion would instead be carried by a saturation acting on the resolved observable-TYPE capacities ---
    a singlet--doublet coincidence of the projection lock A4 (Born--Infeld saturation) on the observable types,
    rather than the shell-alignment (depth) realisation the corpus uses so far.
    That type-saturation arena is not yet constructed, but it is not forbidden; this is the live conceptual frontier
    of the front (see the companion admissibility-coincidence reconnaissance).

  \section{Maximum-entropy interpretation}
    \label{sec:maxent}

    The value $\CV = 1$, i.e.\ $\sigma_r = \bar r$, is the coefficient of variation of the exponential distribution,
    the maximum-entropy distribution on $\mathbb R_+$ at fixed mean.
    The Koide relation can therefore be read as the statement that the square-root masses carry the relative
    dispersion of a maximum-entropy positive variable at fixed mean.
    Unlike a change of carrier, this is a variational principle rather than a numerical coincidence, and the
    projective framework contains exactly one forced entropic object, the projection entropy $\SPi > 0$ of the
    non-injective projection~\cite{Beau2026l}.
    As defined, however, $\SPi = -\sum_{o} \nu(o) \log \nu(o)$ is the Shannon entropy of the pushforward measure
    $\nu$ on the observable space, and the corpus result is its positivity, not a variational extremum over the
    square-root masses~\cite{Beau2026l}; a maximum-entropy \emph{principle} on $\rl$ would have to be supplied, by
    $\SPi$ upgraded to an extremal law over the square-root masses or by a variational spectral entropy acting on
    the same spectral sector.

    \paragraph{The crux is not the entropy but the constraint.}
    The discriminant between derivation and restatement is not the existence of an entropy but the choice of the
    constraint held fixed, because $\CV = 1$ is selected by \emph{one specific} single-constraint problem.
    For a positive variable, the maximum-entropy coefficient of variation depends on which datum is fixed:
    \begin{center}
      \begin{tabular}{lll}
        \hline
        constraint held fixed & maximum-entropy law & $\CV$ \ ($Q = (1+\CV^2)/3$) \\
        \hline
        first moment $\langle \sqrt m \rangle$ (scale) & exponential & $1.000$ \ ($Q = 2/3$, Koide) \\
        second moment $\langle m \rangle$ (norm) & half-normal & $0.756$ \ ($Q \simeq 0.52$) \\
        support $[0,B]$ only & uniform & $0.577$ \ ($Q \simeq 0.44$) \\
        \hline
      \end{tabular}
    \end{center}
    Only the \emph{first-moment} constraint gives $\CV = 1$.
    This is the sharp point: the structurally natural conserved quantity of an amplitude is its quadratic norm
    $\langle m \rangle = \sum_g (\sqrt{m_g})^2$, and fixing \emph{that} yields the half-normal value $\CV = 0.756$,
    not Koide.
    Reaching $\CV = 1$ requires fixing instead the \emph{linear} sum $\langle \sqrt m \rangle = \tfrac13 \sum_g
    \sqrt{m_g}$ and nothing else --- not the variance, not the support.
    Why the mean alone, and why the linear sum rather than the norm, is the entire content of the open test; it is
    not self-evident.

    \paragraph{A structural hook, and where it says not to look.}
    One feature of the projection would make the first-moment constraint natural: the non-injective pushforward
    combines fibre content \emph{linearly}, $\nu(o) = \sum_{\omega \in \Pi^{-1}(o)} \mu(\omega)$~\cite{Beau2026l}.
    If the square-root mass $\sqrt{m_g}$ is identified with such a linearly summed fibre content, then the total
    $\sum_g \sqrt{m_g}$ is the conserved datum and the first-moment (mean-only) constraint is forced, giving
    $\CV = 1$.
    This is the precise, checkable form of the open test, and it is the same identification question that decides
    whether $\SPi$ acts on $\rl$ at all.
    It also fixes where \emph{not} to look, by tension with Proposition~\ref{prop:ladder}: the committed Yukawa
    carrier fixes both the scale $\lambda_Y$ (the mean) and the level dispersion $u$ (the variance), i.e.\ it fixes
    mean \emph{and} variance and lands at $\CV \le 1/\sqrt2$.
    A maximum-entropy object with the mean alone fixed is therefore \emph{not} the cascade-ladder eigenvalue
    vector, whose dispersion is separately pinned; it must be a distinct object --- a linearly summed fibre
    amplitude, scale-fixed and dispersion-free.
    Thus the two statements combine into a definite negative: the maximum-entropy carrier of $\CV = 1$ cannot be
    the square-root of the projected level ladder.

    \paragraph{A lock of the continuous max-entropy route only: the discrete-to-continuous bridge.}
    On the continuous max-entropy route, even granting the first-moment constraint the value $\CV = 1$ does not
    follow for three discrete masses.
    Unit coefficient of variation is a property of the continuous exponential density; three discrete square-root
    masses do not inherit it, and two facts make this precise.
    First, maximising the Shannon entropy of the three weights $p_g = \sqrt{m_g} / \sum_h \sqrt{m_h}$ at fixed
    normalisation gives the uniform distribution $p_g = 1/3$, hence equal square-root masses and $\CV = 0$: the
    degenerate case, the opposite of Koide.
    Second, natural three-point realisations of the maximum-entropy exponential are recipe-dependent and none is
    pinned at unity: its tercile conditional means give $\CV \simeq 0.81$, its $\{1/6, 1/2, 5/6\}$ quantiles give
    $\CV \simeq 0.76$, and the expected sample coefficient of variation of three independent draws is
    $\simeq 0.65$.
    The continuous unit dispersion therefore does not transfer to three values by any canonical discretisation.
    The genuine open question, one level below the constraint question above, is whether there is a forced
    underlying continuous generation coordinate --- exponential by maximum entropy at fixed
    $\langle \sqrt m \rangle$ --- of which the three generations are forced realisations whose discrete coefficient
    of variation is exactly $1$.
    Absent such a forced bridge, unit dispersion is a target invariant, not a derivation.
    This obstruction, however, is specific to the present max-entropy route, in which $\CV = 1$ is read off a
    continuous probability density; it is NOT a universal obstruction to unit dispersion.
    A route that imposes unit dispersion directly on the discrete triple --- for instance the singlet--doublet
    equipartition of the cyclic form~\eqref{eq:cyclic}, where three phases $120^\circ$ apart inherit the
    amplitude-fixed dispersion $\CV = A/\sqrt2$ exactly and independently of the phase $\delta$ (the cyclic form
    being the exact real discrete Fourier transform of any triple) --- does not incur the transfer at all: the
    discrete triple has $\CV = 1 \iff A = \sqrt2$ directly.
    This is exactly how the ternary-phase-carrier note~\cite{Beau2026tpc} addresses the bridge (speculatively),
    relocating the lock to the amplitude (equipartition) condition $A = \sqrt2$ --- which is the condition the
    admissibility-coincidence reading then attacks.
    The discrete-to-continuous bridge is therefore a lock of the continuous max-entropy route only, already
    dissolved on the discrete cyclic route.

    \paragraph{Status of the route.}
    The route is open but now sharply located.
    Its constraint side is the first-moment condition: the projection must fix the linear sum $\sum_g \sqrt{m_g}$
    alone, since the structurally natural quadratic norm gives instead $\CV = 0.756$.
    Its deeper side is the discrete-to-continuous bridge above, which the constraint condition does not close.
    A derivation requires both; if either is chosen by hand, unit dispersion remains a restatement of Koide rather
    than a consequence of the projection.

  \section{Status and open test}
    \label{sec:status}

    This note does not derive the charged-lepton masses.
    Its purpose is to isolate the intrinsic shape constraint carried by the charged-lepton Koide relation and to
    express it as a unit coefficient-of-variation condition on the square-root mass vector.
    The result is a target invariant for the fermionic-mass sector, not a mass formula.

    The status is fourfold.
    The reformulation $Q_\ell = (1 + \CV(\rl)^2)/3$ and the equivalence $Q_\ell = 2/3 \iff \CV(\rl) = 1$ are exact
    (Proposition~\ref{prop:cv}).
    The cascade-ladder bound $\CV \le 1/\sqrt2$ is exact for the pinned ladder (Proposition~\ref{prop:ladder}), but
    it constrains the STATIC LEVELS ($\CV(\rl) \in [0.09, 0.37]$ for their square roots), not the physical masses,
    which carry the depth amplification $A(n_g)$; the physical dispersion is set by the cascade depths, so unit
    dispersion is reached by no presently \emph{constructed} mechanism --- the type-saturation arena on the
    resolved observables (a singlet--doublet coincidence of the A4 projection lock) being not yet built, though not
    forbidden.
    The maximum-entropy reading is an open test, and its crux is the constraint, not the entropy: $\CV = 1$ is
    selected only by a first-moment (mean-only) condition on $\sqrt m$, whereas the natural quadratic norm gives
    $\CV = 0.756$; the route is forced only if the projection fixes the linear sum $\sum_g \sqrt{m_g}$ alone, on a
    linearly summed fibre observable, and this cannot be the level-ladder eigenvalue vector (whose variance is
    separately pinned).
    A deeper lock sits below this constraint question, specific to the max-entropy route: unit dispersion is a
    property of the continuous exponential, and three discrete masses do not inherit it (the discrete-weight
    extremum is uniform, giving $\CV = 0$), so a forced continuous-to-discrete bridge, not merely the constraint, is
    required for a derivation; a route that fixes unit dispersion directly on the discrete triple, as the
    singlet--doublet equipartition of the cyclic form~\eqref{eq:cyclic}, avoids this transfer and relocates the lock
    to the amplitude condition.
    Finally, the concrete hierarchy --- the node placement fixed by the phase $\delta$ --- is deferred and is
    almost certainly the harder problem, since it requires a fine, exponentially sensitive placement of one
    generation against a zero of the square-root profile; it is where the genuine residue of the charged-lepton
    hierarchy remains.
    The clean statement of the easy target --- unit relative dispersion of the square-root masses --- is the
    working conclusion; the mechanism that would force it, and the harder phase problem, are open.

  \newpage

  \bibliographystyle{plain}
  \bibliography{cosmochrony-bibliography}

\end{document}
